\documentclass[12pt,pstricks,border=4pt]{standalone}
\usepackage{amsmath}
\usepackage[dvipsnames]{xcolor}
\usepackage{pst-optexp}

\definecolor{Probe}{HTML}{D99A00}
\definecolor{ProbePale}{HTML}{F7E8BB}
\definecolor{Reference}{HTML}{666D73}
\definecolor{ReferencePale}{HTML}{E7E9EA}
\definecolor{Common}{HTML}{2878B5}
\definecolor{CommonPale}{HTML}{B8D5EA}
\definecolor{Faraday}{HTML}{009E73}
\definecolor{FaradayPale}{HTML}{8FD3BE}
\definecolor{Affected}{HTML}{497D91}
\definecolor{AffectedPale}{HTML}{B8D7DD}
\definecolor{Mixed}{HTML}{4C7F82}
\definecolor{MixedPale}{HTML}{C5DEDC}
\definecolor{PortH}{HTML}{C03C8B}
\definecolor{PortHPale}{HTML}{B7DADB}
\definecolor{PortV}{HTML}{5146A3}
\definecolor{PortVPale}{HTML}{D3D3E0}
\definecolor{Glass}{HTML}{D9EEF3}
\definecolor{Dot}{HTML}{76538B}
\definecolor{DotPale}{HTML}{D9CBE2}
\definecolor{Ink}{HTML}{202428}
\definecolor{SoftInk}{HTML}{60666B}

\psset{linecolor=Ink,linewidth=0.8pt}
\newpsstyle{OptComp}{linecolor=Ink,linewidth=0.9pt,fillstyle=solid,fillcolor=Glass}

\def\AtomsGlyph{%
  \psellipse[fillstyle=solid,fillcolor=black!18,linecolor=Ink,linewidth=0.8pt](0,0)(0.14,0.31)%
}

\def\PhaseDotGlyph{%
  \psline[linecolor=SoftInk,linestyle=dotted,dotsep=1.2pt,linewidth=0.6pt](0,-1.05)(0,1.05)%
  \pscircle[fillstyle=solid,fillcolor=Dot,linecolor=Ink,linewidth=0.7pt](0,0){0.125}%
}

\def\StopGlyph{%
  \psline[linecolor=SoftInk,linestyle=dotted,dotsep=1.2pt,linewidth=0.6pt](0,-1.05)(0,1.05)%
  \pscircle[fillstyle=solid,fillcolor=Ink,linecolor=Ink](0,0){0.12}%
}

\def\BrightCameraGlyph{%
  \psframe[fillstyle=solid,fillcolor=black!20,linecolor=Ink,linewidth=0.9pt](0,0)(1.65,1.55)%
  \psellipse[fillstyle=solid,fillcolor=white,linestyle=none](0.82,0.78)(0.22,0.51)%
}

\def\DarkCameraGlyph{%
  \psframe[fillstyle=solid,fillcolor=Ink,linecolor=Ink,linewidth=0.9pt](0,0)(1.65,1.55)%
  \psellipse[fillstyle=solid,fillcolor=black!22,linestyle=none](0.82,0.78)(0.22,0.51)%
}

\def\FaradayBeforeGlyph{%
  \pscircle[fillstyle=solid,fillcolor=ProbePale,linecolor=Ink,linewidth=0.6pt](0,0){0.49}%
  \psline[linecolor=Probe,linewidth=1.5pt]{<->}(-0.27,0)(0.27,0)%
}

\def\FaradayAfterGlyph{%
  \pscircle[fillstyle=solid,fillcolor=ReferencePale,linecolor=Ink,linewidth=0.6pt](0,0){0.49}%
  \psline[linecolor=Common,linewidth=0.9pt]{<->}(-0.27,0)(0.27,0)%
  \rput{24}{\psline[linecolor=Faraday,linewidth=1.5pt]{<->}(-0.27,0)(0.27,0)}%
}

\def\PBSGlyph{%
  \psframe[fillstyle=solid,fillcolor=Glass,linecolor=Ink,linewidth=0.9pt](-0.39,-0.39)(0.39,0.39)%
  \psline[linecolor=Ink,linewidth=0.8pt](-0.39,0.39)(0.39,-0.39)%
}

\def\MirrorGlyph{%
  \psline[linecolor=Ink,linewidth=2.3pt](-0.50,0.50)(0.50,-0.50)%
}


\begin{document}
\psset{unit=1cm}
\begin{pspicture}(0.15,-2.35)(11.82,4.25)
  \rput[l](0.55,3.45){\textbf{\small DPFI}}

  \pnode(0.55,2.60){Start}
  \pnode(1.50,2.60){AtomNode}
  \pnode(3.80,2.60){LoneNode}
  \pnode(6.80,2.60){PbsNode}
  \pnode(6.80,1.45){FoldNode}
  \pnode(9.20,2.60){LtwoTopNode}
  \pnode(6.80,0.20){LtwoBottomNode}
  \pnode(11.50,2.60){CamTopNode}
  \pnode(6.80,-2.10){CamBottomNode}

  \begin{optexp}
    \optdipole[compname=Atoms,optdipolesize=0.04 1.34,position=end]
      (Start)(AtomNode){\AtomsGlyph}{}
    \lens[compname=Lone,lensheight=2.05,lensradius=2.25,position=end]
      (AtomNode)(LoneNode){}
    \beamsplitter[compname=PBS,bssize=2.90,bsstyle=cube]
      (LoneNode)(PbsNode)(FoldNode){}
    \lens[compname=LtwoTop,lensheight=2.05,lensradius=2.15,position=end]
      (PbsNode)(LtwoTopNode){}
    \lens[compname=LtwoBottom,lensheight=2.05,lensradius=2.15,position=end]
      (FoldNode)(LtwoBottomNode){}

    \backlayer{%
      % Yellow is the linearly polarised probe before the atoms.
      \pspolygon[fillstyle=solid,fillcolor=ProbePale,linestyle=none]
        (0.40,2.00)(1.50,2.00)(1.50,3.20)(0.40,3.20)%
      % Before the PBS, this is the same complete atom-modified transmitted
      % field shown in the PCI and DGI schematics.
      \pspolygon[fillstyle=solid,fillcolor=ProbePale,linestyle=none]
        (1.50,2.00)(3.80,2.00)(5.78,2.37)(5.78,2.83)(3.80,3.20)(1.50,3.20)%
      \pspolygon[fillstyle=solid,fillcolor=AffectedPale,opacity=0.60,linestyle=none]
        (1.50,2.34)(1.50,2.86)(3.80,3.62)(5.78,3.62)%
        (5.78,1.58)(3.80,1.58)%
      % The unaffected probe supplies the same bright background in both
      % +/-45-degree ports.
      \pspolygon[fillstyle=solid,fillcolor=ProbePale,linestyle=none]
        (7.83,2.37)(9.20,2.00)(11.50,2.00)(11.50,3.20)(9.20,3.20)(7.83,2.83)%
      \pspolygon[fillstyle=solid,fillcolor=ProbePale,linestyle=none]
        (6.57,1.58)(6.20,0.20)(6.20,-2.10)(7.40,-2.10)(7.40,0.20)(7.03,1.58)%
      % Transparent overlays show the opposite atom-dependent changes carried
      % on top of those equal backgrounds; they are not separate beams.
      \pspolygon[fillstyle=solid,fillcolor=PortH,opacity=0.20,linestyle=none]
        (7.83,1.58)(9.20,1.58)(11.50,2.34)(11.50,2.86)(9.20,3.62)(7.83,3.62)%
      \pspolygon[fillstyle=solid,fillcolor=PortV,opacity=0.20,linestyle=none]
        (5.78,1.58)(5.78,0.20)(6.54,-2.10)(7.06,-2.10)(7.82,0.20)(7.82,1.58)%
    }
  \end{optexp}

  \psline[linecolor=Ink,linewidth=0.8pt](5.78,3.63)(7.83,1.58)

  % A half-wave plate at 22.5 degrees rotates the analysis basis by 45
  % degrees before the fixed-axis PBS.  The dashed line marks the thin plate
  % without assigning it an artificial longitudinal extent.
  \psline[linecolor=Ink,linestyle=dashed,dash=2.5pt 1.8pt,linewidth=1.05pt]
    (4.85,0.76)(4.85,3.94)
  \rput[r](4.68,0.87){\small half-wave plate, $22.5^{\circ}$}

  % Matched detector planes terminate equal-length, equal-magnification relays.
  \psline[linecolor=Ink,linewidth=1.2pt](11.50,1.75)(11.50,3.45)
  \psline[linecolor=Ink,linewidth=1.2pt](5.95,-2.10)(7.65,-2.10)

  \rput(7.25,3.30){PBS}
  \rput[b](9.20,3.86){$L_{2,H}$}
  \rput[l](7.95,0.20){$L_{2,V}$}
\end{pspicture}
\end{document}
